% --- Archon physics formalization source begin ---
% archon:physics
% archon:covers PhyXMiniProblems/problem_phyx_mini_0379.lean
% archon:source-report reports/phyx_mini/problem_phyx_mini_0379.source.json

\chapter{Physics problem phyx\_mini\_0379}
\label{ch:PhyXMiniProblems_problem_phyx_mini_0379}

\paragraph{Problem source.}
\#\# Physical scenario

The figure shows the thermal energy of 0.14 mol of gas as a function of temperature.

\#\# Question

What is the \textbackslash{}C\_V\textbackslash{} for this gas?

\#\# Answer choices

- A: A: 29\textbackslash{} \textbackslash{}mathrm\{J\textbackslash{}cdot\textbackslash{}mol\textasciicircum{}\{-1\}\textbackslash{}cdot\textbackslash{}K\textasciicircum{}\{-1\}\}

- B: B: 2.97\textbackslash{} \textbackslash{}mathrm\{J\textbackslash{}cdot\textbackslash{}mol\textasciicircum{}\{-1\}\textbackslash{}cdot\textbackslash{}K\textasciicircum{}\{-1\}\}

- C: C: 2.9\textbackslash{} \textbackslash{}mathrm\{J\textbackslash{}cdot\textbackslash{}mol\textasciicircum{}\{-1\}\textbackslash{}cdot\textbackslash{}K\textasciicircum{}\{-1\}\}

- D: D: 129\textbackslash{} \textbackslash{}mathrm\{J\textbackslash{}cdot\textbackslash{}mol\textasciicircum{}\{-1\}\textbackslash{}cdot\textbackslash{}K\textasciicircum{}\{-1\}\}

\#\# Figure caption (auxiliary; use the image as primary evidence)

The image is a graph depicting a linear relationship between two variables: \textbackslash{}( E\_\{\textbackslash{}text\{th\}\} \textbackslash{}) on the y-axis and \textbackslash{}( T \textbackslash{}) on the x-axis. Here are the details:

- **Axes:**
  - The y-axis is labeled \textbackslash{}( E\_\{\textbackslash{}text\{th\}\} \textbackslash{}) (Joules) and has values at 0, 1092, 1492, and 1892.
  - The x-axis is labeled \textbackslash{}( T \textbackslash{}) (°C) and includes values at 0, 100, and 200.

- **Data Representation:**
  - A linear, orange line represents the relationship, starting at approximately (0, 1092) and ending at (200, 1892). 
  - Dashed lines extend from points on the line to the axes to indicate values, such as from the x-axis points 100 and 200 to the line, and from the line to the y-axis at 1492 and 1892.

- **Additional Features:**
  - There is a small break (a jagged line) in the y-axis near the origin, indicating a skipped scale between 0 and 1092.

\#\# Dataset metadata

Specific Heat and Calorimetry; Physical Model Grounding Reasoning, Numerical Reasoning

\paragraph{Recorded answer/context.}
A: A: 29\textbackslash{} \textbackslash{}mathrm\{J\textbackslash{}cdot\textbackslash{}mol\textasciicircum{}\{-1\}\textbackslash{}cdot\textbackslash{}K\textasciicircum{}\{-1\}\}

\paragraph{Figure/image path.}
/root/proposal\_for\_physic/science-mango/phyx\_mini\_run/phyx\_data/test\_image/379.png

\paragraph{Formalization target.}
create a compiling Lean file with sorry bodies at `PhyXMiniProblems/problem\_phyx\_mini\_0379.lean`. The Lean declarations must preserve the physical quantities, dimensions or dimensional roles, figure labels, governing-law hypotheses, and final relation expressed by this problem.
Use Mathlib/Physlib names found through LeanExplore where available. If a physics API is missing, introduce faithful local abstractions rather than scalar placeholder aliases.

\begin{theorem}[Physics formalization target]
\label{thm:physics:phyx_mini_0379:target}
\lean{PhyXMiniProblems.ProblemPhyXMini0379.constantVolumeMolarHeatCapacity_from_thermalEnergyGraph}
\uses{def:physics:phyx-mini-0379:phyxminiproblems-problemphyxmini0379-gasthermalenergyexperiment, def:physics:phyx-mini-0379:phyxminiproblems-problemphyxmini0379-matchesgasthermalenergyscenario, def:physics:phyx-mini-0379:phyxminiproblems-problemphyxmini0379-hasphysicalthermalparameters, def:physics:phyx-mini-0379:phyxminiproblems-problemphyxmini0379-obeysconstantvolumethermalenergylaw, def:physics:phyx-mini-0379:phyxminiproblems-problemphyxmini0379-matchesprimarythermalenergyfigure, def:physics:phyx-mini-0379:phyxminiproblems-problemphyxmini0379-answermolarheatcapacityjoulespermolekelvin}
The assigned autoformalize agent should translate this physics problem into Lean declarations in the covered file, with theorem and lemma proof bodies written as `by sorry`.
\end{theorem}
\begin{proof}
This is an autoformalization task, not a proof task. Produce statements that can later be proved by the physics prover without weakening the physical model.
\end{proof}

% --- Archon declaration topology begin ---
\section{Declaration topology}
\begin{definition}[energy In Joules]
  \label{def:physics:phyx-mini-0379:phyxminiproblems-problemphyxmini0379-energyinjoules}
  \lean{PhyXMiniProblems.ProblemPhyXMini0379.energyInJoules}
  Read a dimensionful thermal energy in SI joules
\end{definition}
\begin{proof}
  This is the declared model definition of energy In Joules.
\end{proof}
\begin{definition}[temperature In Celsius]
  \label{def:physics:phyx-mini-0379:phyxminiproblems-problemphyxmini0379-temperatureincelsius}
  \lean{PhyXMiniProblems.ProblemPhyXMini0379.temperatureInCelsius}
  Read an absolute Physlib temperature in degrees Celsius. Celsius and kelvin increments have the same numerical size, which is used by the difference law below
\end{definition}
\begin{proof}
  This is the declared model definition of temperature In Celsius.
\end{proof}
\begin{definition}[Temperature Mark]
  \label{def:physics:phyx-mini-0379:phyxminiproblems-problemphyxmini0379-temperaturemark}
  \lean{PhyXMiniProblems.ProblemPhyXMini0379.TemperatureMark}
  The three temperature coordinates singled out by the graph
\end{definition}
\begin{proof}
  This is the declared model definition of Temperature Mark.
\end{proof}
\begin{definition}[Axis Label]
  \label{def:physics:phyx-mini-0379:phyxminiproblems-problemphyxmini0379-axislabel}
  \lean{PhyXMiniProblems.ProblemPhyXMini0379.AxisLabel}
  Literal roles of the labels on the horizontal and vertical axes
\end{definition}
\begin{proof}
  This is the declared model definition of Axis Label.
\end{proof}
\begin{definition}[Heat Capacity Kind]
  \label{def:physics:phyx-mini-0379:phyxminiproblems-problemphyxmini0379-heatcapacitykind}
  \lean{PhyXMiniProblems.ProblemPhyXMini0379.HeatCapacityKind}
  The thermodynamic heat-capacity distinction relevant to the question
\end{definition}
\begin{proof}
  This is the declared model definition of Heat Capacity Kind.
\end{proof}
\begin{definition}[Curve Shape]
  \label{def:physics:phyx-mini-0379:phyxminiproblems-problemphyxmini0379-curveshape}
  \lean{PhyXMiniProblems.ProblemPhyXMini0379.CurveShape}
  The shape of the orange data curve drawn in the primary image
\end{definition}
\begin{proof}
  This is the declared model definition of Curve Shape.
\end{proof}
\begin{definition}[Thermal Energy Temperature Graph]
  \label{def:physics:phyx-mini-0379:phyxminiproblems-problemphyxmini0379-thermalenergytemperaturegraph}
  \lean{PhyXMiniProblems.ProblemPhyXMini0379.ThermalEnergyTemperatureGraph}
  \uses{def:physics:phyx-mini-0379:phyxminiproblems-problemphyxmini0379-temperaturemark, def:physics:phyx-mini-0379:phyxminiproblems-problemphyxmini0379-axislabel, def:physics:phyx-mini-0379:phyxminiproblems-problemphyxmini0379-curveshape}
  The plotted curve together with the numerical scale and visible guide marks. The curve itself maps physical absolute temperatures to dimensionful thermal energies; reals in the remaining fields are explicitly labeled axis readouts
\end{definition}
\begin{proof}
  This is the declared model definition of Thermal Energy Temperature Graph.
\end{proof}
\begin{definition}[Gas Thermal Energy Experiment]
  \label{def:physics:phyx-mini-0379:phyxminiproblems-problemphyxmini0379-gasthermalenergyexperiment}
  \lean{PhyXMiniProblems.ProblemPhyXMini0379.GasThermalEnergyExperiment}
  \uses{def:physics:phyx-mini-0379:phyxminiproblems-problemphyxmini0379-heatcapacitykind, def:physics:phyx-mini-0379:phyxminiproblems-problemphyxmini0379-thermalenergytemperaturegraph}
  The gas sample whose molar heat capacity is to be inferred. The two real fields are named numerical readouts in mol and J/(mol K); they are not transparent aliases for physical temperature or energy
\end{definition}
\begin{proof}
  This is the declared model definition of Gas Thermal Energy Experiment.
\end{proof}
\begin{definition}[Matches Gas Thermal Energy Scenario]
  \label{def:physics:phyx-mini-0379:phyxminiproblems-problemphyxmini0379-matchesgasthermalenergyscenario}
  \lean{PhyXMiniProblems.ProblemPhyXMini0379.MatchesGasThermalEnergyScenario}
  \uses{def:physics:phyx-mini-0379:phyxminiproblems-problemphyxmini0379-gasthermalenergyexperiment}
  The prose setup, excluding the graph's measured energy values
\end{definition}
\begin{proof}
  This is the declared model definition of Matches Gas Thermal Energy Scenario.
\end{proof}
\begin{definition}[Has Physical Thermal Parameters]
  \label{def:physics:phyx-mini-0379:phyxminiproblems-problemphyxmini0379-hasphysicalthermalparameters}
  \lean{PhyXMiniProblems.ProblemPhyXMini0379.HasPhysicalThermalParameters}
  \uses{def:physics:phyx-mini-0379:phyxminiproblems-problemphyxmini0379-gasthermalenergyexperiment}
  Positivity and nondegeneracy conditions for the physical branch
\end{definition}
\begin{proof}
  This is the declared model definition of Has Physical Thermal Parameters.
\end{proof}
\begin{definition}[Obeys Constant Volume Thermal Energy Law]
  \label{def:physics:phyx-mini-0379:phyxminiproblems-problemphyxmini0379-obeysconstantvolumethermalenergylaw}
  \lean{PhyXMiniProblems.ProblemPhyXMini0379.ObeysConstantVolumeThermalEnergyLaw}
  \uses{def:physics:phyx-mini-0379:phyxminiproblems-problemphyxmini0379-energyinjoules, def:physics:phyx-mini-0379:phyxminiproblems-problemphyxmini0379-gasthermalenergyexperiment}
  The constant-volume thermal-energy law ΔE\_th = n C\_V ΔT. It is stated for arbitrary pairs of temperatures on the sample's curve and does not assign a numerical value to C\_V
\end{definition}
\begin{proof}
  This is the declared model definition of Obeys Constant Volume Thermal Energy Law.
\end{proof}
\begin{definition}[Matches Primary Thermal Energy Figure]
  \label{def:physics:phyx-mini-0379:phyxminiproblems-problemphyxmini0379-matchesprimarythermalenergyfigure}
  \lean{PhyXMiniProblems.ProblemPhyXMini0379.MatchesPrimaryThermalEnergyFigure}
  \uses{def:physics:phyx-mini-0379:phyxminiproblems-problemphyxmini0379-energyinjoules, def:physics:phyx-mini-0379:phyxminiproblems-problemphyxmini0379-temperatureincelsius, def:physics:phyx-mini-0379:phyxminiproblems-problemphyxmini0379-gasthermalenergyexperiment}
  Exact scalar readouts and visible features transcribed from the primary image. The plotted points are (0,1092), (100,1492), and (200,1892) in (°C,J). No heat-capacity value occurs in this figure-data structure
\end{definition}
\begin{proof}
  This is the declared model definition of Matches Primary Thermal Energy Figure.
\end{proof}
\begin{definition}[Answer Choice]
  \label{def:physics:phyx-mini-0379:phyxminiproblems-problemphyxmini0379-answerchoice}
  \lean{PhyXMiniProblems.ProblemPhyXMini0379.AnswerChoice}
  The four answer labels printed in the source
\end{definition}
\begin{proof}
  This is the declared model definition of Answer Choice.
\end{proof}
\begin{definition}[answer Molar Heat Capacity Joules Per Mole Kelvin]
  \label{def:physics:phyx-mini-0379:phyxminiproblems-problemphyxmini0379-answermolarheatcapacityjoulespermolekelvin}
  \lean{PhyXMiniProblems.ProblemPhyXMini0379.answerMolarHeatCapacityJoulesPerMoleKelvin}
  \uses{def:physics:phyx-mini-0379:phyxminiproblems-problemphyxmini0379-answerchoice}
  Numerical value beside each answer, in joules per mole-kelvin
\end{definition}
\begin{proof}
  This is the declared model definition of answer Molar Heat Capacity Joules Per Mole Kelvin.
\end{proof}
\begin{definition}[recorded Answer Choice]
  \label{def:physics:phyx-mini-0379:phyxminiproblems-problemphyxmini0379-recordedanswerchoice}
  \lean{PhyXMiniProblems.ProblemPhyXMini0379.recordedAnswerChoice}
  \uses{def:physics:phyx-mini-0379:phyxminiproblems-problemphyxmini0379-answerchoice}
  The dataset metadata records answer choice A
\end{definition}
\begin{proof}
  This is the declared model definition of recorded Answer Choice.
\end{proof}
% --- Archon declaration topology end ---
% --- Archon physics formalization source end ---
